% =====================================================================
% Informe — Pregunta 2: Mapeo de inundaciones con Sentinel-1 y U-Net
% Curso: Redes Neuronales y Aprendizaje Profundo — Maestría en IA, UNI
% Plantilla: IEEEtran (conference). Objetivo: 8 páginas (rango permitido 6–10).
% Las figuras van en la carpeta figuras/ del proyecto de Overleaf.
% =====================================================================
\documentclass[conference]{IEEEtran}
\usepackage[utf8]{inputenc}
\usepackage[T1]{fontenc}
\usepackage[spanish,es-noshorthands,es-tabla]{babel}
\usepackage{amsmath,amssymb}
\usepackage{graphicx}
\usepackage{booktabs}
\usepackage{multirow}
\usepackage{siunitx}
\usepackage{xcolor}
\usepackage[hidelinks]{hyperref}
\usepackage{cite}
\graphicspath{{figuras/}}
\sisetup{output-decimal-marker={.}, separate-uncertainty=true}

% Marca de pendiente: se borra cuando la sección esté escrita
\newcommand{\pendiente}[1]{\textcolor{red}{[Pendiente: #1]}}

% \seccion{archivo}{contenido provisional}: si existe secciones/archivo.tex lo incluye;
% si todavía no existe, muestra el contenido provisional. Así cada sección nueva solo se
% sube como archivo, sin volver a tocar este main.tex.
\newcommand{\seccion}[2]{\IfFileExists{secciones/#1.tex}{\input{secciones/#1}}{#2}}

\begin{document}

\title{Mapeo de inundaciones con Sentinel-1 y una U-Net:\\
ablaciones, generalización a Lima y fusión con Sentinel-2}

\author{%
\IEEEauthorblockN{Andhersson M. Salazar F.\IEEEauthorrefmark{1},
Alexander A. Chavez C.\IEEEauthorrefmark{2},
Alex R. Pomazón S.\IEEEauthorrefmark{3} y
Christian A. Urviola G.\IEEEauthorrefmark{4}}
\IEEEauthorblockA{Maestría de Especialización en Inteligencia Artificial\\
Universidad Nacional de Ingeniería, Lima, Perú\\
\IEEEauthorrefmark{1}andhersson.salazar.f@uni.pe, ORCID: 0009-0006-4517-5666\\
\IEEEauthorrefmark{2}alexander.chavez.c@uni.pe, ORCID: 0009-0000-0956-9082\\
\IEEEauthorrefmark{3}alex.pomazon.s@uni.pe, ORCID: 0009-0004-8514-1546\\
\IEEEauthorrefmark{4}christian.urviola.g@uni.pe, ORCID: 0000-0003-4621-7128}}

\maketitle

% ---------------------------------------------------------------------
\begin{abstract}
El radar de apertura sintética (SAR) de Sentinel-1 permite mapear inundaciones bajo nubes,
pero su principio físico (el agua calma se ve oscura) tiene excepciones. Entrenamos una U-Net
implementada desde cero con las polarizaciones VV y VH sobre Sen1Floods11 y la evaluamos con un
protocolo de tres semillas. La pérdida combinada BCE + Dice y las conexiones de salto son las
decisiones más influyentes, mientras que las etiquetas débiles de Otsu no mejoran al modelo. En
test el modelo alcanza un IoU de agua de 0.669, frente a 0.554 de Otsu, y se mantiene en el
evento no visto de Bolivia (0.663). En cambio, colapsa en la inundación de Lima de marzo de 2017
(IoU 0.003). Con imágenes antes y después del evento mostramos que allí el agua urbana y de
huaico \emph{se aclara} en el radar. Un análisis de fallos confirma con pruebas cuantitativas
cuatro mecanismos (incidencia rasante, doble rebote, agua no especular y superficies secas
lisas) y descarta un quinto. Finalmente, fusionar seis bandas de Sentinel-2 eleva el IoU de test
a 0.848, aunque la ganancia proviene casi por completo del óptico y está favorecida por cómo se
etiquetaron los datos.
\end{abstract}

\begin{IEEEkeywords}
SAR, Sentinel-1, segmentación semántica, U-Net, inundaciones, Sen1Floods11, generalización geográfica
\end{IEEEkeywords}

% ---------------------------------------------------------------------
% Extensión final: ~10 páginas con figuras (límite permitido: 10).
% ---------------------------------------------------------------------

\seccion{01_introduccion}{%
\section{Introducción}
\label{sec:introduccion}
\pendiente{motivación: Niño Costero 2017 en Lima; el SAR ve a través de nubes y de noche;
el agua calma se ve oscura (reflexión especular); objetivo y contribuciones (E1--E6)}}

\seccion{02_relacionado}{%
\section{Trabajo relacionado}
\label{sec:relacionado}
\pendiente{U-Net y CNN sobre S1 frente a umbrales; conjuntos de datos; inundación urbana;
fusión S1+S2; modelos fundacionales; carencias que motivan este trabajo}}

\seccion{03_datos}{%
\section{Datos}
\label{sec:datos}
\pendiente{Sen1Floods11 (dorados, débiles, particiones); Lima: S1 del 17/03/2017 y EMSN038;
preprocesamiento y desbalance (9.4\,\% de agua)}}

\seccion{04_metodologia}{%
\section{Metodología}
\label{sec:metodologia}
\pendiente{arquitectura (E1), pérdidas, métricas agregadas, entrenamiento, protocolo
(3 semillas, selección en valid, test intacto)}}

\section{Experimentos y resultados}
\label{sec:resultados}

\seccion{05a_ablacion_perdida}{%
\subsection{Ablación (b): función de pérdida y conexiones de salto}
\pendiente{E2(b)}}

\seccion{05b_ablacion_debiles}{%
\subsection{Ablación (a): etiquetas débiles}
\pendiente{E2(a)}}

\seccion{05c_test}{%
\subsection{Evaluación en distribución (test)}
\pendiente{E3}}

\seccion{05d_generalizacion}{%
\subsection{Generalización geográfica: Bolivia y Lima}
\pendiente{E4}}

\seccion{05e_fallos}{%
\subsection{Análisis cualitativo de fallos}
\pendiente{E5}}

\seccion{05f_fusion}{%
\subsection{Fusión con Sentinel-2 (bono)}
\pendiente{E6}}

\seccion{06_discusion}{%
\section{Discusión crítica y limitaciones}
\label{sec:discusion}
\pendiente{sesgo de las etiquetas hacia el óptico; Lima: etiqueta de huella máxima vs.\ SAR;
una sola fecha; tamaño de muestra; decisiones no exploradas}}

\seccion{07_conclusiones}{%
\section{Conclusiones}
\label{sec:conclusiones}
\pendiente{3--4 conclusiones y trabajo futuro (detección de cambios, multitemporal)}}

\section*{Disponibilidad del código}
El código (módulos \texttt{src/}, cuadernos por fase y resultados en CSV) está disponible en el
repositorio público del proyecto en GitHub.
% Enlace al repositorio: reemplazar USUARIO/REPOSITORIO y quitar el % de la línea siguiente.
\url{https://github.com/AndherssonSalazar/mapeo-inundaciones-sar-unet.git}

\bibliographystyle{IEEEtran}
\bibliography{referencias}

\end{document}